\section{引言}

数据流通是指跨组织、跨平台和跨计算环境的受控数据使用。在这类场景中，可信并不只意味着存储文件是否保持不变。一个可部署系统还必须回答数据使用是否经过授权、计算过程是否可审计、证据是否具有防篡改性，以及事故发生后责任是否能够定位。金融、医疗、政务和工业协作场景通常避免直接传输明文数据，而是依赖隐私保护计算、可信执行、数据连接器或受控 API。因此，校验对象从静态文件扩展为数据分片、计算日志、模型更新、任务摘要、结果哈希和恢复事件。

经典云存储审计主要回答远程数据是否仍被正确保存。可证明数据持有（PDP）通过挑战响应协议和同态标签降低校验成本~\cite{ateniese_provable_2007}。可恢复性证明（PoR）进一步强调数据在有限损坏后仍应保持可恢复~\cite{juels_pors_2007,shacham_compact_2013}。动态 PDP 支持经过认证的插入、删除和修改~\cite{erway_dynamic_2015}。这些技术为远程完整性校验提供了密码学基础，但审计过程通常由用户或第三方审计者发起。在多方数据流通工作流中，审计者自身可能延迟、遗漏或错误报告审计结果，因此审计轨迹也必须能够被独立验证。

基于区块链的审计将问题重心从“数据是否可校验”转向“审计过程本身是否可信”。已有工作使用区块哈希或随机数作为不可预测挑战来源，利用智能合约记录审计和惩罚，并通过联盟链共识约束拖延或恶意审计者~\cite{xue_identity-based_2019,zhang_blockchain-based_2021,fan_dredas_2020,lin_consortium_2022}。在同一方向上，BLoP 融合区块链、可信计算和隐私保护计算~\cite{Zhang2025BloPAT}。它通过远程度量和监测组件维护可信节点，检测异常事件，并在联盟链场景下使用 PBFT 风格确认机制替代工作量证明。这种全节点确认模型适用于强一致性和高保障要求优先于成本约束的场景。

本文研究的问题更聚焦：在将 BLoP 作为全节点可信审计基线的前提下，能否为频繁出现的低风险或中风险审计事件增加一条轻量化路径，同时保留回退到全节点确认的能力。该问题重要的原因在于，真实数据流通工作流可能产生大量常规审计记录。如果每条记录都由所有可信节点确认并逐条上链，系统可能将大量资源消耗在常规证据处理上，而非高风险事件。TrustAuditFlow 通过将审计工作流划分为两条路径来处理这一问题：常规事件通过 Merkle 根和哈希链锚定进行批量提交，并由从 BLoP 风格可信节点集合中选出的可信委员会确认；关键事件或争议事件则回退到全节点路径。

核心研究挑战并不只是减少区块链写入次数。该设计还必须保留可验证证据链，避免委员会成为新的单点信任，并将完整性失败与恢复过程连接起来。因此，TrustAuditFlow 将区块链视为时间线和问责层，而不是重型证明计算或大规模证据存储的位置。数据分片、审计证明和恢复材料保留在链下；链上记录承诺、摘要、恢复状态和信任分变化。

本文贡献如下：

\begin{itemize}
  \item 面向高频审计场景，提出 BLoP 全节点审计模型的一种轻量化扩展，同时保留全节点确认作为高保障回退路径。
  \item 定义面向数据流通的闭环审计恢复流程：分片存储、抽样校验、异常检测、纠删码辅助恢复、链上锚定和信任分更新。
  \item 设计信任分驱动的审计委员会机制，在不将委员会替代完整可信节点集合的前提下降低常规确认成本。
  \item 给出分层架构，使昂贵的校验和证据存储留在链下，而 Merkle 根、哈希链锚点、恢复状态和问责记录提交到链上。
  \item 刻画轻量化路径的安全边界和成本因素，包括委员会选择风险、批量提交成本、回退条件，以及与 PDP/PoR 模块的关系。
\end{itemize}

